% Generated draft: requires mathematical and source review.
\documentclass[11pt]{article}
\usepackage[margin=25mm]{geometry}
\usepackage{fontspec}
\setmainfont{Latin Modern Roman}
\usepackage{amsmath,amssymb,mathtools,unicode-math}
\setmathfont{Latin Modern Math}
\usepackage{enumitem,needspace,xurl,hyperref}
\setlength{\parindent}{0pt}
\setlength{\parskip}{5pt}
\setlength{\emergencystretch}{3em}
\newcommand{\N}{\mathbb{N}}
\newcommand{\Z}{\mathbb{Z}}
\newcommand{\Q}{\mathbb{Q}}
\newcommand{\R}{\mathbb{R}}
\newcommand{\C}{\mathbb{C}}
\newcommand{\F}{\mathbb{F}}
\newcommand{\A}{\mathbb{A}}
\newcommand{\PP}{\mathbb{P}}
\newcommand{\E}{\mathbb{E}}
\newcommand{\eps}{\varepsilon}
\newcommand{\dd}{\,\mathrm{d}}
\begin{document}
\section*{UnsolvedMath ID: 30005407}
\textbf{Source problem code:} catalogue entry
\subsection{Definitions and mathematical statement}
The supplied record concerns a particular ring \(S\) and ideal \(I\subseteq S\), whose definitions and hypotheses are absent from the supplied excerpt. They must be recovered from the cited report; the question must not be interpreted as quantified over all rings and ideals.

For a positive integer \(d\), let \(I^d\) denote the ordinary ideal power. Write \(\operatorname{Ext}^3_S(M,S)\) for the third right derived functor of \(\operatorname{Hom}_S(M,S)\), and \(\ell_S(M)\) for the composition length of an \(S\)-module \(M\), when finite. Provided the indicated lengths are finite for all sufficiently large positive integers \(d\), put
\[
 a_d=\frac{\ell_S\!\left(\operatorname{Ext}^3_S(S/I^d,S)\right)}{d^6}.
\]
The source question is whether
\[
 \lim_{\substack{d\to\infty\\ d\in\Z_{>0}}}a_d
\]
exists for the intended \(S\) and \(I\). The supplied excerpt does not establish the required finiteness hypotheses or specify whether existence means a finite real limit or allows an extended real limit. Consequently this is an incomplete formulation pending recovery of the original setup.

The excerpt identifies this question as an instance of higher \(\varepsilon\)-multiplicities, citing its reference [5]. It separately reports that, when \(\operatorname{char}(k)=2\), calculations in its reference [7, Theorem 1.9] imply that
\[
 \lim_{\substack{d\to\infty\\ d\in\Z_{>0}}}
 \frac{\ell_S\!\left(\operatorname{Ext}^3_S(I^d/I^{d+1},S)\right)}{d^5}
\]
does not exist. This is supplied contextual information about a different sequence, not a conclusion about \(a_d\). The field \(k\), its relationship to \(S\), and the characteristic restrictions on the target question are not specified. The passage also describes explicit intermediate cohomology calculations for line bundles on \(X\), or sheaves rendered as \(D_dR(e)\) on \(P\), as open in general and related to determinantal ideals; these objects are undefined in the supplied excerpt.
\subsection{Short English statement}
For the specific ring and ideal intended in the report, ask whether the lengths of \(\operatorname{Ext}^3_S(S/I^d,S)\), divided by \(d^6\), have a limit as positive integers \(d\) tend to infinity. The setup and length finiteness assumptions are missing. The supplied record labels the question open and relates it to higher \(\varepsilon\)-multiplicities; these are source qualifications requiring review.
\subsection{Sources}
\begin{itemize}
\item Local UnsolvedMath registry, supplied record with permanent ID \(30005407\), source problem code \texttt{catalogue entry}, title: \emph{Normalized Ext Limits and Epsilon Multiplicity}. No individual proposer is supplied; the proposed year is 2023.
\item \emph{Resolutions in Local Algebra and Singularity Theory} (2023), Oberwolfach Reports, pp.\ 327--395, DOI \texttt{10.4171/owr/2023/6}: \url{https://doi.org/10.4171/owr/2023/6}. The extracted passage carries page number 343. This report documents workshop discussions and is also the supplied literature source labeled ``original OWR report.''
\item The excerpt cites [5] for higher \(\varepsilon\)-multiplicities and [7, Theorem 1.9] for the related characteristic-two calculation. Their bibliographic entries are not supplied.
\item Literature assessment checked at \texttt{2026-08-22}: the supplied curation labels the problem open and reports that no primary source found by its cutoff proves existence in the stated generality. This is supplied dated curation, not an independently verified literature assessment.
\item The registry labels the statement \texttt{corrected\_verified} and reports that the official report was checked directly and definitions restored. These are supplied verification claims; the present draft has not independently checked the report, and the supplied record still omits the definitions of \(S\) and \(I\).
\end{itemize}
\end{document}
